\section{Driver layer}

\subsection{Interfaces and state objects}

Hardware-independent interfaces are defined in \file{src/bdx_slow_control/drivers/base.py}. State is returned with immutable dataclasses such as \code{PowerChannelState}, \code{ChillerState}, and \code{DaqCrateState}.

Current abstract interfaces include:

\begin{itemize}
  \item \code{PowerSupplyDriver};
  \item \code{HighVoltageDriver};
  \item \code{ChillerDriver};
  \item \code{SensorDriver};
  \item \code{DaqCrateDriver}.
\end{itemize}

All drivers implement \code{ping()}. Read methods return complete state objects. Write methods perform one explicit device action and raise an exception when the operation cannot be completed or verified according to the driver contract.

\subsection{Simulation first}

A subsystem should have a simulation backend before or together with its hardware implementation. Simulation enables:

\begin{itemize}
  \item IOC and display development without equipment;
  \item deterministic validation of writable PV behavior;
  \item fault and interlock testing;
  \item continuous tests on machines without the hardware network.
\end{itemize}

Set \code{driver.simulation=true} for simulation implementations so \pv{SIMULATION} and \pv{COMM_STATUS} are correct.

\subsection{Adding a hardware backend}

Use the following sequence.

\begin{enumerate}
  \item Identify the existing abstract interface and state dataclass. Extend them only when the new device requires information that is meaningful at the subsystem contract level.
  \item Add the implementation under \file{src/bdx_slow_control/drivers/hardware/}. Keep socket, serial, protocol parsing, locking, timeouts, and device command strings inside this module.
  \item Provide a small builder function that validates required JSON fields and returns the driver instance.
  \item Register the driver name in \file{src/bdx_slow_control/drivers/factory.py}. Unsupported names must fail with an explicit configuration error.
  \item Add a JSON example or operational profile entry. Do not embed host addresses or channel counts in IOC code.
  \item Add unit tests for command formatting, parsing, limit behavior, timeouts, and representative failures. Use fake transports or local test servers where possible.
  \item Run the subsystem IOC against simulation, then perform read-only hardware checks, then controlled write tests under an approved procedure.
\end{enumerate}

\begin{importantbox}[Event-loop behavior]
A slow blocking hardware call must not stall unrelated caproto Channel Access traffic. Use the established driver/IOC pattern for offloading and serialization when a device library or network protocol is blocking.
\end{importantbox}

\section{IOC layer}

\subsection{ManagedIOC}

\file{src/bdx_slow_control/iocs/common.py} defines \code{ManagedIOC}. A subsystem IOC should normally:

\begin{itemize}
  \item declare PVs with caproto \code{pvproperty};
  \item implement \code{poll_device()} to write readbacks from one coherent driver state;
  \item implement putters for explicit writes;
  \item catch write failures, call \code{mark_failure()}, and re-raise or return a rejected status as appropriate;
  \item avoid transport-specific code and deployment constants.
\end{itemize}

\subsection{Staged operator requests}

Routine hardware changes should prefer request/apply semantics. The low-voltage IOC is the reference pattern:

\begin{enumerate}
  \item request PVs are initialized from the first successful device read;
  \item the operator edits voltage and current requests without immediate hardware writes;
  \item \pv{APPLY_CMD} validates both values together;
  \item the driver applies the request;
  \item the IOC rereads the device and updates readbacks;
  \item \pv{APPLY_STATUS} and \pv{APPLY_MESSAGE} record the outcome.
\end{enumerate}

If a multi-step hardware write can partially succeed, the failure message must state that possibility and the IOC should reread the device before reporting the final state.

\subsection{Commands}

Momentary command PVs should normally return to false after execution. Confirmed actions in Phoebus are required for operations such as output enable, chiller start, all-off, simulated interlock injection, and interlock reset.

Global or device all-off callbacks are registered in \code{PrototypeContext}. A reset clears the latched global interlock state but must not automatically restore outputs.

\section{Builders and CLI entry points}

\file{src/bdx_slow_control/builders.py} converts validated JSON objects into IOC groups and PV databases. Driver selection occurs through the factory. \file{src/bdx_slow_control/prototype.py} combines available subsystem files, checks for duplicate PVs, and enforces compatible server interfaces.

Console scripts are declared in \file{pyproject.toml}. The principal development commands are:

\begin{lstlisting}[style=shell]
bdx-prototype-ioc --config-dir config/profiles/default
bdx-psu-ioc --config config/profiles/default/psu.json
bdx-chiller-ioc --config config/profiles/default/chiller.json
bdx-environment-ioc --config config/profiles/raspberry/environment.json
bdx-pv-list --config-dir config/profiles/default
bdx-generate-displays --config-dir config/profiles/default \
  --output-dir phoebus/displays
\end{lstlisting}

\subsection{Adding a new subsystem}

A new subsystem normally requires all of the following:

\begin{enumerate}
  \item a state model and abstract driver interface in the driver layer;
  \item a simulation implementation;
  \item one or more IOC groups under \file{src/bdx_slow_control/iocs/};
  \item a builder function in \file{builders.py};
  \item a driver-factory function or registration;
  \item a key in the \code{BUILDERS} mapping;
  \item an individual CLI entry point when focused execution is useful;
  \item a console-script entry in \file{pyproject.toml};
  \item one or more JSON profiles;
  \item generated Phoebus coverage and tests;
  \item Archiver PV-list updates when the subsystem must be archived;
  \item documentation updates.
\end{enumerate}

The subsystem filename used by the aggregated prototype must match its \code{BUILDERS} key, for example \code{psu} maps to \file{psu.json}.
